\documentclass{exam}

\usepackage{amsmath,amssymb,amsfonts}
\usepackage{graphicx}
\usepackage{amsthm}
\usepackage[french]{babel}
\usepackage{enumitem}

\setlength{\topmargin}{-.5in} \setlength{\textheight}{9.25in}
\setlength{\oddsidemargin}{0in} \setlength{\textwidth}{6.8in}

\makeatletter
\def\gap#1{%
  \setbox\@tempboxa\hbox{{\color{white}#1}}%
  \@tempdima\wd\@tempboxa
  \rlap{\unhbox\@tempboxa}%
  \hb@xt@\@tempdima{\dotfill}%
}
\makeatother   

\theoremstyle{definition}
\newtheorem{exercice}{Exercice}

\begin{document}
\Large

\noindent{\bf Contrôle 1 -- 1h. \hfill Terminales -- Math expertes \hfill 02/10/2024}

\medskip\hrule
\vspace{2em}

\begin{exercice}Exprimer les nombres complexes suivants sous leur forme algébrique:
    \begin{enumerate}
        \item $3-(2+3i)(4+i)$,
        \item $\frac{3+i}{2-i}$.
    \end{enumerate}
\end{exercice}

\vspace{1em}

\begin{exercice}\ 
    
    \begin{enumerate}
        \item Résoudre l'équation  $z + 2i \overline{z} = 1 - i$ dans $\mathbb{C}$.
        \item Résoudre l'équation $z^2 - \overline{z} + 1 = 0$ dans $\mathbb{C}$.
        \item Décrire l'ensemble des solutions complexes de $\overline{z} + iz = 0$. À quelle figure géométrique correspond-t-il ?
    \end{enumerate}
    
\end{exercice}

\vspace{1em}

\begin{exercice}
    Soit $Z= \frac{z+i}{z-i}$ pour tout $z\neq i$.
    \begin{enumerate}
        \item Exprimer $\overline{Z}$.
        \item En déduire tous les nombres complexes $z$ tels que $Z$ soit réel.
    \end{enumerate}
\end{exercice}

\vspace{1em}

\begin{exercice}
    On considère la fonction $f$ qui a tout $z \in \mathbb{C}$ associe $f(z) = z^2+2z+9$.
    \begin{enumerate}
        \item Calculer l'image de $-1+i\sqrt{3}$ par $f$.
        \item On pose $z = x + iy$, donner l'écriture sous forme algébrique de $f(z)$.
        \item Quels sont les nombres complexes qui possèdent un image réelle par $f$ ?
    \end{enumerate}
\end{exercice}

\vspace{1em}

\begin{exercice}\ 
    \begin{enumerate}
        \item Donner la forme algébrique de $(2+3i)^4$ à l'aide du binôme de newton.
        \item Écrire une formule inspirée du binôme de newton pour $(a-b)^n$.
        \item En déduire que $\sum_{k=0}^{n} (-1)^k \binom{n}{k} = 0$.
    \end{enumerate}
\end{exercice}

\vspace{1em}

\begin{exercice}
    On considère la suite $u_n = i^n$ définie pour tout $n \in \mathbb{N}$.
    \begin{enumerate}
        \item Quelle est la nature de cette suite ?
        \item En déduire la forme algébrique de $z = \sum_{k=0}^{2021} i^k$.
    \end{enumerate}
\end{exercice}

\end{document}